\section{Descripcion Extensa de la Actividad técnica (10-20 paginas)}

\subsection{Analisis y optimizacion inicial de rendimiento en VR}
El primer encargo tecnico estuvo centrado en mejorar la experiencia en Realidad Virtual desde el punto de vista del rendimiento percibido. Se revisaron elementos de escena, configuraciones de render y decisiones de pipeline que afectaban a los FPS y a la estabilidad temporal. Se hizo uso extensivo del las herramientas de diagnostico de Unity, como el \emph{Profiler} y el \emph{Frame Debugger}, para identificar cuellos de botella en CPU y GPU, y se aplicaron optimizaciones como la reduccion de \textit{draw calls}, la simplificacion de materiales y la gestion más eficiente de las jerarquias de objetos. Se observó que los materiales utilizados podían ser optimizados para permitir \textit{batching}. Esto se logró mediante la unificación de materiales y la reducción de variantes, lo que permitió que Unity agrupara objetos con el mismo material en una sola llamada de dibujo, reduciendo significativamente los draw calls y mejorando el rendimiento general.\par

En esta fase se ajustaron configuraciones de build y paquetes para alinear el proyecto con el objetivo VR, se revisaron escenas y prefabs XR para reducir problemás visuales y mejorar el confort del usuario. Entre otros, llegó un nuevo escenario por parte del equipo de arte. Toda la escena principal, de la clase, fue portada al nuevo escenario, asegurando que los elementos de interacción y de animación se mantuvieran funcionales.\par

Posteriormente, se comprobo la compatibilidad entre sistemás de entrada y rastreo de movimientos en VR, incluyendo la coexistencia de componentes legacy y del New Input System cuando fue necesario para mantener el funcionamiento en distintas rutas de ejecucion. Debido a esta disparidad, se unificó el uso de los sistemás de input, migrando exclusivamente al New Input System. El bloque se cerro con una validacion del flujo de build y de la configuracion XR orientada al dispositivo objetivo.\par

Este trabajo inicial permitio establecer una base más estable para el resto de desarrollos funcionales, reduciendo incidencias de tirones y errores de integracion entre editor y build final.\par

\subsection{Diseno y programación de estados para describir comportamientos}
Una vez estabilizado el arranque técnico del proyecto, se inicio el trabajo de diseno de una herramienta para modelar comportamientos del alumnado mediante grafos y maquinas de estados. Esta linea resultaba especialmente relevante para la evolucion hacia comportamientos más complejos, incluyendo casos de alumnado no normativo. Se redacto el documento inicial de diseno técnico del sistema de grafo de comportamiento. Junto a ello se revisaron las transiciones de estado y sus condiciones, introduciendo criterios más estrictos de validacion para detectar inconsistencias durante la ejecucion. Con estas decisiones quedo preparada la base para extender de forma modular nuevos conflictos y sus respuestas animadas.\par

Como resultado, el proyecto paso de una configuracion menos estructurada a un enfoque más mantenible, con mejor trazabilidad entre conflicto recibido, transicion de estado aplicada y animacion reproducida.\par

El diseño de esta herramienta se pospuso parcialmente para poder avanzar con otras tareas de mayor prioridad. Sin embargo, ahora, cerca de la finalización de este ``corte vertical'' de la aplicación, es posible que se retome esta línea para completar la implementación de la herramienta de diseño de comportamientos, pues la serie de comportamientos que se plantean en toda su complejidad quizá requiera de una refactorización significativa del código que maneja los estados y las transiciones. En este sentido, este apartado podría considerarse un punto de partida para una evolución futura que permita consolidar esta parte del proyecto con una herramienta de diseño visual que facilite la incorporación de nuevos casos de uso pedagógicos sin necesidad de cambios extensos en el código base.\par

\subsection{Robustecimiento de conexion web, mensajeria y aserciones}
El bloque con mayor continuidad temporal fue la verificacion y mejora de la conexion con la web a la que se enlaza la aplicacion VR. El objetivo principal fue aumentar la fiabilidad y la capacidad de diagnostico ante mensajes incompletos, inesperados o no implementados. Se reforzo el sistema de logging y de errores para facilitar la depuracion de mensajes y de rutas de ejecucion no implementadas. Tambien se introdujeron comprobaciones explicitas para eventos web desconocidos y aserciones sobre condiciones de transicion, distinguiendo entre modos más estrictos o más relajados segun el contexto y la dirección del tutor de prácticas. En paralelo, se ajusto la creacion de botones web para garantizar la correspondencia correcta entre interfaz y conflicto objetivo, y se integró un selector de servidor que permitiera alternar entre entornos de desarrollo y producción y otros futuros.\par

Parte del sistema fue reestructurada con una orientacion más dirigida a datos, mediante descriptores y pipelines de generación de conflictos. El objetivo era que la incorporacion de nuevos tipos no exigiera reescrituras extensas de codigo procedural, sino ampliaciones localizadas y más seguras. Con esta base se redujo la fragilidad del flujo de conexion, mejoró la observabilidad de los fallos y se habilito una operativa de pruebas más realista para el resto del equipo al poder alternar servidores de forma controlada.\par

En esto último, cuando hubo que expandir la herramienta para soportar cinco nuevos conflictos asociados a TEA y TDAH, el propio pipeline de generación de conflictos fue guiando al desarrollador a través de los pasos necesarios para asegurar que cada nuevo conflicto se integrara correctamente en el sistema, pues pasaba por validaciones explícitas en cada etapa, desde la definición del descriptor de conflicto, hasta su incorporación en la lógica de generación y su representación en la interfaz. Esto permitió que la ampliación funcional se realizara de forma ordenada y con un menor riesgo de introducir errores o inconsistencias en el sistema existente.\par

\subsection{Reescritura parcial y orientacion a datos para escalabilidad}
Para sostener el aumento de complejidad funcional, se acometieron cambios estructurales en sistemás de gestion de clase y estudiantes. El objetivo fue permitir el escalado de conflictos y simplificar la incorporacion de nuevos tipos sin reescrituras másivas.\par

Se introdujeron los descriptores de conflicto y su representación interna como unión discriminada, se corrigieron problemas de alineacion y de punteros en layouts de datos que podian provocar lecturas incorrectas, y se gestionaron con mayor seguridad los casos limite, como estudiantes insuficientes, repeticion de alumno conflictivo o conflictos no implementados. Tambien se ajusto el pipeline de generación para asegurar que todos los conflictos definidos se procesaran y reportaran correctamente.\par

\subsection{Ampliacion y validacion de nuevos conflictos}
% TODO: rev
Una de las aportaciones nucleares de la practica fue la validacion de los conflictos existentes y la ampliacion funcional con nuevos conflictos asociados a TEA y TDAH. La ampliacion se materializo en la incorporacion de nuevos descriptores de conflicto y en el paso de estos por el pipeline de generación. A nivel de comportamiento, se implementaron nuevas transiciones de maquina de estados para cubrir situaciones no normativas y se ajustaron los gestores de clase y alumnado para asegurar que todos los tipos de conflicto pudieran asignarse y ejecutarse correctamente. Tambien se revisaron estados no implementados y rutas de contingencia para evitar bloqueos de ejecucion.\par

El trabajo no se limito a anadir casos; incluyo una validacion sistematica de compatibilidad con sistemás ya existentes, como la asignacion de alumnos, las condiciones de entrada y salida de estado, y la sincronizacion con la interfaz y la animacion. A nivel funcional, este bloque incremento la riqueza de escenarios disponibles y dejo el sistema preparado para futuras ampliaciones con menor friccion técnica.\par

\subsection{Integracion de animaciones y adaptacion de prefabs}
La incorporacion de nuevos conflictos requirio sincronizacion con contenido audiovisual y con las jerarquias de personajes. Se ajustaron y depuraron los controladores de animacion para los nuevos estados y transiciones, se verificaron los disparadores ligados a las condiciones de conflicto y se adaptaron los prefabs de alumnado y de elementos relacionados para asegurar referencias y componentes correctos en la jerarquia. Tambien se corrigieron comportamientos de sentado, levantado y orientacion en escena.

Con ello se consiguio una representación más consistente del comportamiento en escena, evitando desincronizaciones entre el estado logico y la respuesta visual.

\subsection{Soporte a multiples servidores de prueba}
Para facilitar el trabajo de QA y la coordinacion con otros miembros del equipo, se habilito la seleccion de distintos servidores de conexion desde interfaz. Se anadio un componente de \textit{dropdown} de servidor en el menu y se conecto su seleccion a la logica de conexion para alternar entre entornos, por ejemplo desarrollo y producción. Posteriormente se verifico su funcionamiento en build para evitar dependencias exclusivas del modo editor.

Esta mejora redujo tiempos de prueba y errores de configuracion manual al cambiar de entorno backend.

\subsection{Verificacion y exportacion satisfactoria a Meta Quest 3}
Finalmente, se completaron ajustes de compatibilidad y configuracion para exportar y ejecutar correctamente en Meta Quest 3 sin incidencias graves reportadas. Se revisaron las configuraciones de XR y de entrada relevantes para el dispositivo, se corrigieron camaras y componentes de seguimiento para mantener un trazado estable y se realizaron comprobaciones en las escenas de menu y aula para validar el ciclo de uso principal.

La exportacion funcional a Quest 3 se considera un hito de cierre tecnico de la practica por su impacto en la validacion de producto real.

\subsection{Incidencias técnicas destacadas y resolucion}
Durante el periodo de practicas aparecieron incidencias que requirieron diagnostico especifico. Las más relevantes fueron desajustes entre el conflicto solicitado desde la web y el conflicto realmente activado, estados de animacion no disparados pese a una transicion logica correcta, errores por rutas no implementadas o mensajes inesperados, y casos limite en la asignacion de estudiantes conflictivos. La estrategia general de resolucion se baso en instrumentacion mediante logs y aserciones, validacion incremental y pruebas dirigidas por caso de uso.

\subsection{Valor tecnico aportado al proyecto}
La contribucion realizada no se limito a cambios puntuales, sino que mejoró la base del proyecto en terminos de robustez, extensibilidad y capacidad de prueba. En conjunto, el trabajo dejo un sistema más tolerante a errores y más observable en la capa de conexion y estados, unas estructuras más escalables para el crecimiento de conflictos y comportamientos, y una integracion funcional más completa entre el backend de eventos, la logica de estado y la representación animada en VR.